\section{Introduction}\label{sec:intro}

Many succinct arguments reduce the correctness of a computation to a claim about a multilinear polynomial evaluated at a random point.
This is the case for sumcheck-based systems such as Spartan~\cite{Set20} and HyperPlonk~\cite{CBBZ23}.
Such systems need a \emph{multilinear polynomial commitment scheme}: the prover commits to a multilinear polynomial $\tilde f$ in $n$ variables and later proves that $\tilde f(z) = v$ for a point $z$ chosen by the verifier.
When the scheme is built from hash functions and error-correcting codes, the resulting arguments are transparent and rely on no number-theoretic assumption, which makes them candidates for post-quantum security.

A multilinear polynomial can be given in two forms.
In \emph{table form} (also called evaluation or Lagrange form) it is given by its values $f(b)$ on the Boolean hypercube $B_n = \{0,1\}^n$; this is the form in which provers naturally hold their witnesses.
In \emph{coefficient form} it is given by its monomial coefficients.
The two forms are related by M\"obius inversion on the Boolean lattice.

Code-based schemes in the BaseFold paradigm~\cite{ZCF24} commit to an encoding of $f$ and prove an evaluation by running a sumcheck whose round challenges also drive an FRI-like folding of the codeword~\cite{BBHR18}.
With Reed--Solomon codes on smooth multiplicative subgroups, as in DeepFold~\cite{GLHQTZ25}, the classical FRI fold $U \mapsto U_e + \theta\, U_o$ acts on coefficient vectors.
It computes the multilinear evaluation of the \emph{coefficient} vector: the monomial coefficients of the committed univariate polynomial must be the monomial coefficients of $\tilde f$.
A prover that holds a table must therefore first move to coefficient form, and must keep the two views consistent across the committed word, the sumcheck state and the folds.

This paper asks whether there is a basis of univariate polynomials in which FRI-style folding over a smooth multiplicative domain acts \emph{directly} on the table form.
There is, and it is explicit.

\subsection{Contributions}\label{sec:intro:contrib}

\paragraph{The Boolean-kernel basis (\cref{sec:kernel}).}
For $y \in \Fq^n$ let $K_y(X) = \prod_{k=1}^{n}(X^{2^{k-1}} + y_k)$.
We show that the $N = 2^n$ polynomials $K_b$, $b \in B_n$, which have coefficients in $\{0,1\}$, form a basis of $\F[X]_{<N}$ over every field $\F$.
The change of basis to monomials is the Kronecker power $C^{\otimes n}$ of $C = \bigl(\begin{smallmatrix} 0 & 1 \\ 1 & 1\end{smallmatrix}\bigr)$.
It costs exactly $\tfrac{n}{2}N$ additions or subtractions in either direction, and no multiplications.
In this basis, having degree $< 2^t$ is a sign condition on the coordinates (\cref{thm:lowdeg}).
For all but at most $N-1$ points $\zeta$, evaluating $U$ at $\zeta$ equals, up to an explicit nonzero scalar, evaluating the multilinear extension of its coordinates at a point on an explicit curve (\cref{prop:eval}).

\paragraph{The fold dictionary (\cref{sec:fold}).}
Our main structural result is \cref{thm:fold-dictionary}.
Let $U = \sum_b \lambda(b) K_b$ and $U = U_e(X^2) + X U_o(X^2)$. Then the kernel coordinates of the fold
\[
  \mathrm{pfold}_r(U) \;=\; (2r-1)\,U_e + (1-r)\,U_o
\]
are exactly $(1-r)\lambda(2b') + r\lambda(2b'+1)$, the restriction of the first variable of the table $\lambda$.
Consequently, encoding a table $f$ as the Reed--Solomon codeword of $U_f = \sum_b f(b) K_b$ and folding it $n$ times with challenges $r_1,\dots,r_n$ yields the constant $\tilde f(r_1,\dots,r_n)$ (\cref{cor:codeword-fold}), with no change of representation at any point.
The kernel fold agrees with the classical FRI fold up to a nonzero scalar and a reparametrisation of the challenge (\cref{sec:fold:fri}); mixing the two parametrisations breaks correctness (\cref{ex:pitfall}), a mistake we made in an early implementation.

\paragraph{A multilinear PCS with a self-contained soundness proof (\cref{sec:pcs,sec:sound}).}
We instantiate the BaseFold paradigm with the kernel encoding and prove it sound in the unique-decoding regime (\cref{thm:sound}).
The only external ingredient about codes in the soundness proof is the correlated-agreement theorem for Reed--Solomon codes of Ben-Sasson, Carmon, Ishai, Kopparty and Saraf~\cite{BCIKS23}; the efficiency of the extractor also uses polynomial-time unique decoding of Reed--Solomon codes.
The theorem enters through a one-line observation: the kernel fold of any word is an affine line $u^0 + r\,u^1$ whose generators determine the word (\cref{lem:line}).
The proof tracks the set of passing query chains level by level, and it yields evaluation binding (\cref{cor:binding}).
The same argument gives round-by-round knowledge soundness with an explicit decoding extractor (\cref{thm:rbr}).
The BCS theorem in the quantum random oracle model of Chiesa, Manohar and Spooner~\cite{CMS19} does not cover instances containing a Merkle-committed oracle, such as the commitment. Chiesa, Di, Hu and Zheng~\cite{CDHZ25} extend the BCS transformation to interactive oracle reductions, whose instances may contain such oracles, under a relaxed round-by-round knowledge soundness hypothesis. We derive this hypothesis from \cref{thm:rbr} (\cref{thm:rrbr}); their theorem then makes the non-interactive scheme knowledge sound, for a single opening, against quantum adversaries that choose the commitment, with explicit concrete bounds (\cref{thm:cdhz,cor:pq}); \cref{rem:pq-params} gives parameters for this setting.

\paragraph{Implementation and experiments (\cref{sec:impl,sec:exp}).}
We give a Rust implementation over the Goldilocks field with a quadratic extension for challenges.
We compare it with a coefficient-form baseline that shares all other code.
The two encodings are indistinguishable in running time at the resolution of our machine, as the operation counts predict.
The kernel encoding thus provides its structural property at no measurable cost.

\paragraph{What we do not claim.}
The scheme is not asymptotically or concretely faster than Reed--Solomon instantiations of BaseFold or DeepFold (\cref{rem:basefold}).
Its evaluation protocol is BaseFold's sumcheck-plus-folding, applied to a different encoding.
Our soundness analysis is restricted to the unique-decoding regime, which needs more queries than list-decoding analyses~\cite{Hab24,Zei26}.
Zero knowledge is out of scope, and we analyse one commitment opened at one point.
The post-quantum statement (\cref{cor:pq}) uses the theorem of~\cite{CDHZ25}, which we quote and do not reprove, and it applies to their compiler (salted Merkle trees, their challenge derivation), not to our implementation as it stands (\cref{rem:fs}).
The contribution is the basis and the dictionary: a precise, elementary account of why FRI folding and table-form multilinear restriction are the same operation over multiplicative domains in odd characteristic.

\subsection{Related work}\label{sec:intro:related}

\paragraph{Proximity tests for Reed--Solomon codes.}
FRI~\cite{BBHR18} tests proximity to Reed--Solomon codes by repeated even--odd folding.
Its soundness in the unique-decoding and list-decoding regimes rests on proximity gaps for Reed--Solomon codes~\cite{BCIKS23}.
WHIR~\cite{ACFY24} tests proximity to constrained Reed--Solomon codes with very fast verification and supports multilinear evaluation claims.
Our soundness proof uses only the unique-decoding correlated-agreement theorem of~\cite{BCIKS23}; extending it to the list-decoding regime is the most natural next step.

\paragraph{Multilinear commitments from folding.}
BaseFold~\cite{ZCF24} introduced foldable codes and the interleaving of sumcheck rounds with folding rounds under shared challenges.
Our evaluation protocol follows it directly.
DeepFold~\cite{GLHQTZ25} instantiates the paradigm with Reed--Solomon codes and improves proof size; its folding computes coefficient-form evaluations, and our coefficient-form baseline of \cref{sec:exp} follows the same principle.
Hab\"ock~\cite{Hab24} analyses BaseFold over Reed--Solomon codes up to the Johnson bound.
Zeilberger~\cite{Zei26} proves proximity gaps for multilinear evaluation for all linear codes.
Both results concern the soundness of the folding paradigm; our contribution is orthogonal and concerns the encoding.

\paragraph{Binary fields and the novel polynomial basis.}
FRI-Binius~\cite{DP24} adapts BaseFold to binary towers using the novel polynomial basis of Lin, Chung and Han~\cite{LCH14}.
In that basis the additive FFT, and with it FRI-style folding over additive subspaces, interacts cleanly with multilinear structure.
The Boolean-kernel basis plays an analogous role for smooth \emph{multiplicative} subgroups in \emph{odd} characteristic.
The two constructions are genuinely different. In characteristic $2$ every factor $X^{2^{k-1}} + 1$ equals $(X+1)^{2^{k-1}}$, so $K_b = X^{N-1-b}(X+1)^{b}$ is a Bernstein-type basis; moreover the multiplicative group of a field of characteristic $2$ has odd order, so smooth multiplicative domains do not exist, and the even--odd formulas~\eqref{eq:even-odd} divide by $2$.

\paragraph{Multilinear-to-univariate reductions.}
A different line of work proves multilinear evaluations with a \emph{univariate} polynomial commitment, typically KZG.
Gemini~\cite{BCHO22} reduces a coefficient-form multilinear claim to univariate claims through recursive even--odd splitting.
Zeromorph~\cite{KT24} works in table form and uses quotients and degree checks.
Mercury~\cite{EG25} reaches constant proof size and expresses the evaluation as a coefficient of a product with a tensor-structured polynomial. Samaritan~\cite{GPS25} improves on Gemini in proof size, prover cost and verifier cost.
Mohamed~\cite{Moh25} relates univariate and multilinear sumchecks through the inverse Kronecker substitution.
These schemes treat the univariate commitment as a black box.
In ours, the univariate polynomial is committed through a Reed--Solomon code, and the basis is chosen so that the code's own folding carries the multilinear structure.

\paragraph{Other domains.}
Brakedown~\cite{GLSTW23} and related schemes avoid FFTs by using linear-time encodable codes.
Circle STARKs~\cite{HLP24} obtain smooth folding domains over the Mersenne prime $2^{31}-1$ from the circle curve.
The fold dictionary as stated needs a smooth multiplicative subgroup in odd characteristic.
Whether a Boolean-kernel-type basis exists for circle domains and for additive domains is an open question that we leave to future work.

\subsection{Organisation}\label{sec:intro:org}

\Cref{sec:prelim} fixes notation and collects the facts on multilinear polynomials and Reed--Solomon codes that we use, including the correlated-agreement theorem.
\Cref{sec:proofs} recalls interactive oracle proofs and their probability space, round-by-round (knowledge) soundness, the sumcheck protocol, polynomial commitments, and the definitions and the theorem of~\cite{CDHZ25} on the BCS transformation for interactive oracle reductions, which we use for the non-interactive scheme.
\Cref{sec:kernel} develops the Boolean-kernel basis, and \cref{sec:fold} proves the fold dictionary.
\Cref{sec:pcs} defines the commitment scheme and \cref{sec:sound} proves its soundness.
\Cref{sec:impl,sec:exp} describe the implementation and the measurements, and \cref{sec:concl} lists open problems.
\Cref{app:notation} collects the notation.
